\begin{exercisechunk}
\item \textbf{Construct a two-round induction lookup.}
\label[exercise]{ex:l14-induction-head}
\exerciseprofile{2}{2}{2}{\exproof\quad\exconstruction\quad\exinterpretation}{%
Show how one attention round constructs predecessor--current records and a
second round uses those records for content-based retrieval.}
An induction head retrieves the token that followed an earlier occurrence of
the current token. The construction below expresses that mechanism as two
causal lookup rounds and makes the dependency between the rounds explicit.

Let \(x_1,\ldots,x_n\in\cS\), where \(n\ge3\). Assume that there is a unique
index \(j\in[n-2]\) such that \(x_j=x_n\), and that no other index in
\([n-1]\) contains \(x_n\). Initially, position \(i\) stores \(x_i\) and its
position index. A round may use a fixed positional query or an equality query,
retrieve one causal source state, and combine the retrieved value with the
destination's residual state.
\begin{enumerate}[(a),beginpenalty=10000]
\item Give a first causal round after which every position \(i\ge2\) stores
the record \((x_{i-1},x_i)\). State the source selected at position \(i\) and
the residual information retained there.
\item Give a second causal round at position \(n\) that uses \(x_n\) as a key
and returns \(x_{j+1}\). Identify the unique selected position.
\item Prove that the two-round program returns the token following the unique
earlier occurrence of \(x_n\) for every sequence satisfying the assumptions.
\item Explain why the required key--value record is absent from the round-zero
states. Under the stated one-source-per-round abstraction, identify the two
dependent retrievals that prevent the same construction from being performed
in one round.
\end{enumerate}
\end{exercisechunk}
